% !TeX program = xelatex
% !TeX encoding = UTF-8 Unicode
% !BIB program = biber

% essa classe aceita as seguintes opções:
% abnd ou ieee pra estilo de citaçõ
% english ou french pra língua dos slides (não passar nada é português)
% minted pra usar o pacote minted (configurado para python e C)
\documentclass[minted]{relatorio}

% Todo comando \add trabalha com uma lista e todo comando \set trabalha com um
% único elemento, ou seja, pode haver apenas um \setorientador mas vários
% \addauthor. Todo comando \add e \set tem um \get correspondente (por exemplo,
% \getauthors, \getsubject). Esses comandos são definidos na classe
% relatorio.cls. No caso dos \set, os \get são implicitos, e não você não vai
% encontrá-los na classe.

\title{Exemplo de Relatório}
\addauthor{Álan Crístoffer e Sousa}
\addauthor{Lanna Gontijo Ferreira}
\setsubject{Controle Digital e Laboratório}

\preamble{}

% Essa classe usa o biblatex com backend biber ao invés do bibtex.
\addbibresource{bibliothek.bib}

% Controla o texto das captions das figuras. Neste caso, centraliza.
\captionsetup{justification=centering}

\begin{document}
% !TeX root = document.tex
% !TeX encoding = UTF-8 Unicode

\IEEEtitleabstractindextext{%
\begin{abstract}
  O controle discreto no tempo se aproxima mais da realidade atual, onde o
  controle é feito de forma digital. As leis de controle implementadas
  utilizando sistemas discretos podem ser escritas com funções simples,
  normalmente de soma e multiplicação de termos, que são fáceis de interpretar
  e, principalmente, de implementar em um sistema digital. Por tratar o sistema
  como digital, ele leva em conta problemas da natureza desses sistemas, como o
  tempo de amostragem, normalmente apenas ignorados em implementações contínuas
  nesses sistemas. Este trabalho visa a implementação de dois controladores
  digitais em um sistema massa-mola em plano inclinado, visando demonstrar a
  síntese e implementação de tais controladores, e a análise desses sistemas.
\end{abstract}

\begin{IEEEkeywords}
  controle digital, jury, espaço de estados
\end{IEEEkeywords}
}


\maketitle{}

% !TeX root = document.tex
% !TeX encoding = UTF-8 Unicode

\section{Floats}
\subsection{Figuras}

\begin{slide}{Figura Única}
  \begin{figure}[ht!] % ht! é o posicionamento da figura. Inútil em slides.
    \centering % centraliza a figura no ambiente figure
    % height=0.3\textheight faz a imagem ter 30% da altura do texto. Outra
    % opção seria width=0.5\textwidth, por exemplo. Use a primeira pra
    % figuras formato retrato e a segunda pra figuras formato paisagem.
    % Ajuste o tamanho da figura usando a porcentagem.
    \includegraphics[height=0.3\textheight]{imgs/cefet}
    \caption{Logo do CEFET.} % texto que aparece abaixo da figura.
  \end{figure}
\end{slide}

\begin{slide}{Subfiguras}
  \begin{figure}[ht!]
    \centering
    \subfloat[One]{\includegraphics[height=0.3\textheight]{imgs/cefet}}\qquad
    \subfloat[Two]{\includegraphics[height=0.3\textheight]{imgs/cefet}}
    \caption{Duas logos do CEFET.}
  \end{figure}
\end{slide}

\subsection{Tabela}

\begin{slide}{Tabela}
  % Eu uso https://www.tablesgenerator.com pra gerar a tabela, e depois edito.
  \begin{table}
    \begin{tabular}{l|c|c|c|c} % alinhamento das células: Left/Center/Right
      \toprule
      Professor & Mecânica       & Eletrônica     & Controle       & Programação    \\ \midrule
      Lúcio     & \(\checkmark\) &                &                &                \\
      Valter    &                &                & \(\checkmark\) &                \\
      Thiago    &                &                &                & \(\checkmark\) \\
      Marlon    &                & \(\checkmark\) &                &                \\
      \bottomrule
    \end{tabular}
    \caption{Eixo de professores.}
  \end{table}
  % Este exemplo mostra como usar top/mid/bottom rules e linhas verticais. No
  % entanto, o povo do latex não recomenda usar isso dizendo que fica feio. No
  % fim do dia o documento é seu, use ou não se quiser.
\end{slide}

% !TeX root = document.tex
% !TeX encoding = UTF-8 Unicode

\section{Section 2}
\subsection{Subsection 2}

\begin{slide}{List}
    \vspace*{\fill}
    \begin{enumerate}
        \item One
        \item Two
        \item Three
    \end{enumerate}
    \vspace*{\fill}
\end{slide}

% !TeX root = document.tex
% !TeX encoding = UTF-8 Unicode

\section{Equações}
\subsection{Numeradas}

\begin{slide}{Numeradas}
  Equation:
  \begin{equation}
    \dot{x} = Ax+Bu
  \end{equation}

  \vspace*{\fill} % Adiciona espaço de forma que tudo fique igualmente espaçado

  Align:
  \begin{align}
    \dot{x} & = Ax+Bu \\
    y       & =Cx+Du
  \end{align}

  \vspace*{\fill}

  Equation+Aligned:
  \begin{equation}
    \begin{aligned}
      \dot{x} & = Ax+Bu \\
      y       & =Cx+Du
    \end{aligned}
  \end{equation}
\end{slide}

\subsection{Não Numeradas}

\begin{slide}{Não Numeradas}
  Equation:
  \begin{equation*}
    \dot{x} = Ax+Bu
  \end{equation*}

  \vspace*{\fill} % Adiciona espaço de forma que tudo fique igualmente espaçado

  Align:
  \begin{align*}
    \dot{x} & = Ax+Bu \\
    y       & =Cx+Du
  \end{align*}

  \vspace*{\fill}

  Equation+Aligned:
  \begin{equation*}
    \begin{aligned}
      \dot{x} & = Ax+Bu \\
      y       & =Cx+Du
    \end{aligned}
  \end{equation*}
\end{slide}

\begin{slide}{Align Com algumas numeradas}
  \begin{align}
    \dot{x} & = Ax+Bu           \\
    y       & = Cx+Du           \\
    c       & = Ex+Fu \nonumber % remove a numeração de uma equação
  \end{align}
\end{slide}

% !TeX root = document.tex
% !TeX encoding = UTF-8 Unicode

\section{Citações}
\subsection{Citações}

Para referenciar, há basicamente duas formas: \texttt{parencite} e
\texttt{textcite}. Eles são os equivalentes do biblatex aos comandos
\texttt{citep} e \texttt{citet} do bibtex. O primeiro serve para citar no fim do
parágrafo:

\begin{quote}
  Há quem diga que a lua é um queijo~\parencite{book:ogata}.
\end{quote}

O segundo para citar quando o nome do autor faz parte do texto:

\begin{quote}
  Segundo~\textcite{book:ogata}, a lua é um queijo.
\end{quote}

% !TeX root = document.tex
% !TeX encoding = UTF-8 Unicode

\section{Códigos}
\subsection{Códigos}

O pacote \texttt{minted} oferece um ambiente para apresentação de código com
\textit{syntax highlight}.

Python 3:

\begin{python3code}
#!/usr/bin/env python3


def fib(n):
    a = b = 1
    while n > 1:
        a, b, n = b, a + b, n - 1
    return b

import sys

if len(sys.argv) > 1:
    print(fib(int(sys.argv[1])))
else:
    print("This application finds the n'th fibonacci number.\nType the desired value for n:")
    print(fib(int(input())))
\end{python3code}

C:

\begin{ccode}
#include <stdint.h>
#include <stdlib.h>

extern "C" {
    uint64_t fib(uint64_t n)
    {
        uint64_t a = 1;
        uint64_t b = 1;
        uint64_t c = 1;

        while (n-- > 2) {
            c = a + b;
            a = b;
            b = c;
        }

        return b;
    }
}
\end{ccode}


\nocite{*}
\printbibliography{}
\end{document}
